\subsection{Удары}
%%%%%%%%%%%%%%%%%%%%%%%%%%%%%%%%%%%%%%%%%%%%%%%%%%%%
%%%%%%%%%%%%%%%%%%%%%%%%%%%%%%%%%%%%%%%%%%%%%%%%%%%%
%%%%%%%%%%%%%%%%%%%%%%%%%%%%%%%%%%%%%%%%%%%%%%%%%%%%
\z[{\begin{wrapfigure}[3]{r}{0.3\textwidth}%
		\vspace{-10pt}%
		\begin{flushright}%
 			\includegraphics[width=0.3\textwidth]{6.3.1}%
		\end{flushright}%
	\end{wrapfigure}%
	}]%
{На гладком столе лежит брусок массой $2m$. К бруску прикреплена пружинка жёсткости $k$, другой конец которой закреплён к стене. В брусок летит пластилиновый шарик массы $m$ со скоростью $v$. Чему будет равна скорость бруска после удара?}
%
{$u = \frac{1}{3} v$}

%%%%%%%%%%%%%%%%%%%%%%%%%%%%%%%%%%%%%%%%%%%%%%%%%%%%
\z{На сколько максимально сожмётся пружинка в предыдущей задаче?}
%
{$\Delta x = v \sqrt{\frac{m}{3k}}$}

%%%%%%%%%%%%%%%%%%%%%%%%%%%%%%%%%%%%%%%%%%%%%%%%%%%%
\z[{\begin{wrapfigure}[6]{r}{0.2\textwidth}%
		\vspace{-10pt}%
		\begin{flushright}%
 			\includegraphics[width=0.2\textwidth]{6.3.3}%
		\end{flushright}%
	\end{wrapfigure}%
	}]%
{На столе лежит брусок. На лёгкой нити длиной $L$ висит шарик, касаясь бруска. Нить вертикальна. Масса бруска в $7$ раз больше массы шарика. Шарик отклоняют в сторону так, что нить занимает горизонтальное положение, и отпускают. После неупругого удара о брусок шарик останавливается. Найти скорость бруска сразу после удара.}
%
{$v = \sqrt{\frac{2 g l}{7}}$}

%%%%%%%%%%%%%%%%%%%%%%%%%%%%%%%%%%%%%%%%%%%%%%%%%%%%
\z{В предыдущей задаче брусок проехал по горизонтальной поверхности расстояние $S$. Найти коэффициент трения скольжения между бруском и столом.}
%
{$\mu = \frac{L}{49 S}$}

%%%%%%%%%%%%%%%%%%%%%%%%%%%%%%%%%%%%%%%%%%%%%%%%%%%%
\z{Пуля массы $m$, имеющая начальную скорость $v$, пробивает подвешенный на нити свинцовый шар массы $M$ и вылетает из него со скоростью $v/2$. Какая скорость будет у шарика после вылета пули?}
%
{$u = v \cdot \frac{m}{2M}$}

%%%%%%%%%%%%%%%%%%%%%%%%%%%%%%%%%%%%%%%%%%%%%%%%%%%%
\z{В предыдущей задаче найдите, какая часть кинетической энергии пули перешла в тепло.}
%
{$\eta = \frac{3}{4} - \frac{m}{4M}$}

%%%%%%%%%%%%%%%%%%%%%%%%%%%%%%%%%%%%%%%%%%%%%%%%%%%%
\z{Два тела одинаковой массы летели во взаимно перпендикулярных направлениях с одинаковой по модулю скоростью. Произошло абсолютно неупругое столкновение. Какая часть кинетической энергии перешла в тепло?}
%
{Половина}


